\scanpage{8}
\mspage{1}
1. $K$ complexe semi-simplicial

$\unsure{\mathscr P|P}$ \deleted{Catégorie} $\otimes$ Catégorie AU

représentation \inserted{$\Lambda$} de $K$ dans $\mathscr P$
\[
\text{données}\left\{
\begin{aligned}
&\text{a) }\forall u\in K_1,\qquad\Lambda(u)\in\mathscr P\\
&\text{b) }\forall\sigma\in K_2,\ \text{posant }u=d_1\sigma,\ v=d_3(\sigma),\ w=d_{\unsure{2}}(\sigma)\\
&\hspace{9em}\Lambda(u)\otimes\Lambda(w)\simeq\Lambda(\unsure{v|u|w})
\end{aligned}\right.
\]
axiome \quad c) $\forall\Delta\in K_3$, on a ``axiome du pentagone

La composée \quad $K_1\longrightarrow\mathscr P\longrightarrow\pi_0(\mathscr P)$ définit un ``homo\uline{morph}isme'' \quad $K_1\xrightarrow{\alpha}\pi_0(\mathscr P)=G$

d'où un ``homomorphisme'' \quad $K_1\xrightarrow{\bar\alpha}\operatorname{Aut}\big(\underbrace{\pi_1(\mathscr P)}_{M}\big)$

d'où un système local sur $K$ de valeur $\underline M$, noté $(\unsure{{}_\alpha M|M},\bar\alpha)=\bar M$

\uline{Pb} Quand $K,\mathscr P,\deleted{\unsure{(M,\bar\alpha)}}\inserted{\alpha\ \text{(donc }\bar M\text{)}}$ sont fixés, on veut ``classer'' les représentations de $K$ \unsure{dans|ds} $\mathscr P$ correspondantes.

% slanted
\margin{$\operatorname{REP}(K;\mathscr P,\alpha)$\\$\operatorname{REP}(\mathcal X;\mathscr P,\alpha)$}
\uline{Résultat} Ces représentations forment une 2-catégorie, qui est un ``2-torseur'' sous une \uline{2-catégorie de Picard} \inserted{stricte} d'invariants
\[
H^2(K,\unsure{(}\bar M\deleted{\unsure{,\bar\alpha)|,\alpha)}}),\quad
H^1(K,\unsure{(}\bar M\deleted{\unsure{,\bar\alpha)|,\alpha)}}),\quad
H^0(K,\unsure{(}\bar M\deleted{\unsure{,\bar\alpha)|,\alpha)}}).
\]
Cette 2-catégorie est contravariante p.r. à $K$ variable, \deleted{\illegible{2}} covariante p.r. à $\mathscr P$ variable

Cas particulier utilisé par la suite : $K$ est le \uline{nerf} d'une catégorie $\mathcal X$. \deleted{\unsure{On}}

2. $\mathscr P$ Gr-catégorie \deleted{\unsure{de Picard}} donnée, d'invariants
\[
\pi_0(\mathscr P)=\unsure{C|G}\ ;\quad\pi_1(\mathscr P)=M
\]
On va lui associer une 2-gerbe \inserted{$\widetilde{\mathscr P}$} sur $B_G$, donc à tout $X\in\operatorname{ob}B_G$ (ensemble où $\unsure{G|C}$ opère) on va associer une catégorie $\widetilde{\mathscr P}(X)$, ainsi :

Soit $\mathcal X$ le nerf de la catégorie \inserted{$\mathcal X$} sous-jacente à $X$
% archive sheet number: 7
